
\subsection{拉格朗日交的应用}
\label{subsec:6-1-3}

在构造谱序列的过程中，
我们将证明如下弗洛尔上同调结构定理。

\begin{theorem}{6.1.18}\label{thm:6-1-18}
设 $L$ 是相对 spin 拉格朗日子流形。
则对每个
$(\boldsymbol b,b)\in
\mathcal M_{\mathrm{weak,def}}(L)$，
存在数 $\ell_p,m_p,\lambda_{p,i}$，
使
\begin{equation}
\begin{aligned}
HF^p((L,(\boldsymbol b,b));
\Lambda_{0,\mathrm{nov}})
\cong{}&
\left(
\bigoplus_{i=1}^{m_p}
\frac{\Lambda_{0,\mathrm{nov}}^{(0)}}
{F^{\lambda_{p,i}}
\Lambda_{0,\mathrm{nov}}^{(0)}}
\right)\\
&\quad\oplus
(\Lambda_{0,\mathrm{nov}}^{(0)})^{\oplus\ell_p}
\end{aligned}
\tag{6.1.19.1}\label{eq:6-1-19-1}
\end{equation}
作为 $\Lambda_{0,\mathrm{nov}}^{(0)}$-模的同构。
\end{theorem}

注意，(6.1.19.1) 蕴含
\begin{equation}
HF^p((L,(\boldsymbol b,b));
\Lambda_{\mathrm{nov}})
\cong
(\Lambda_{\mathrm{nov}}^{(0)})^{\oplus\ell_p}
\tag{6.1.19.2}\label{eq:6-1-19-2}
\end{equation}
作为 $\Lambda_{\mathrm{nov}}^{(0)}$-模的同构；
这些模具有 $\mathbb Z$-分次。
回忆
$\Lambda_{0,\mathrm{nov}}
\cong\Lambda_{0,\mathrm{nov}}^{(0)}[e,e^{-1}]$，
而 $e$ 的次数为二。
因此，
$HF^p((L,(\boldsymbol b,b));
\Lambda_{0,\mathrm{nov}})$
作为 $\Lambda_{0,\mathrm{nov}}^{(0)}$-模
以二为周期；
$HF^p((L,(\boldsymbol b,b));
\Lambda_{\mathrm{nov}})$
作为 $\Lambda_{\mathrm{nov}}^{(0)}$-模亦然。

\begin{theorem}{6.1.20}\label{thm:6-1-20}
设 $(L^{(0)},L^{(1)})$
是 $M$ 中一对干净相交的
相对 spin 拉格朗日子流形。
对每个
$\mathbf b_i\in
\mathcal M_{\mathrm{weak,def}}(L^{(i)})$
满足
$\mathfrak{PO}(\mathbf b_0)
=\mathfrak{PO}(\mathbf b_1)$ 及
$\pi_{\mathrm{amb}}(\mathbf b_1)
=\pi_{\mathrm{amb}}(\mathbf b_0)$，
都存在数 $\ell_p,m_p,\lambda_{p,i}$，
使
\begin{equation}
\begin{aligned}
&HF^p((L^{(1)},\mathbf b_1),
(L^{(0)},\mathbf b_0);
\Lambda_{0,\mathrm{nov}})\\
&\qquad\cong
\left(
\bigoplus_{i=1}^{m_p}
\frac{\Lambda_{0,\mathrm{nov}}^{(0)}}
{F^{\lambda_{p,i}}
\Lambda_{0,\mathrm{nov}}^{(0)}}
\right)
\oplus
(\Lambda_{0,\mathrm{nov}}^{(0)})^{\oplus\ell_p}
\end{aligned}
\tag{6.1.21.1}\label{eq:6-1-21-1}
\end{equation}
作为 $\Lambda_{0,\mathrm{nov}}^{(0)}$-模的同构。
\end{theorem}

同样，(6.1.21.1) 蕴含
\begin{equation}
HF^p((L^{(1)},\mathbf b_1),
(L^{(0)},\mathbf b_0);
\Lambda_{\mathrm{nov}})
\cong
(\Lambda_{\mathrm{nov}}^{(0)})^{\oplus\ell_p}
\tag{6.1.21.2}\label{eq:6-1-21-2}
\end{equation}
作为 $\Lambda_{\mathrm{nov}}^{(0)}$-模的同构；
这些模具有 $\mathbb Z$-分次。
仍有
$\Lambda_{0,\mathrm{nov}}
\cong\Lambda_{0,\mathrm{nov}}^{(0)}[e,e^{-1}]$，
且 $\deg e=2$。
所以，
$HF^p((L^{(1)},\mathbf b_1),
(L^{(0)},\mathbf b_0);
\Lambda_{\mathrm{nov}})$
作为 $\Lambda_{\mathrm{nov}}^{(0)}$-模
以二为周期。

把奇、偶两种次数的挠指数合并如下：
\begin{equation}
\{\lambda_i\mid i=1,2,\ldots\}
=
\{\lambda_{p,i}\mid
p=0,1;\ i=1,2,\ldots\}.
\tag{6.1.22}\label{eq:6-1-22}
\end{equation}
若按 $\lambda_i\le\lambda_{i+1}$ 排列，
则计入重数后的各 $\lambda_i$
唯一确定。
称它们为定理 6.1.20 中
拉格朗日对子弗洛尔上同调的
挠指数\liftIndexRangeStart{I217}。
注意，这里使用有理系数；
“挠”指 $\Lambda_{0,\mathrm{nov}}$-模
的挠部分\liftIndex{I016}。
$\ell_0+\ell_1$ 称为
弗洛尔上同调的贝蒂数。

本章最后一节研究拉格朗日交。
其中一个问题是估计
$L$ 与 $\psi(L)$ 的交点数，
这里 $\psi$ 是哈密顿微分同胚。
我们将在第 6.5.1、6.5.2 小节
分别证明定理 H、I。

再次强调：
对单个拉格朗日子流形 $L$，
其弗洛尔上同调
$HF((L,b);\Lambda_{0,\mathrm{nov}})$
与定义过程中所作的各种选择无关。
由定理 4.1.3、4.1.4 还可知，
其“自由部分”的秩也是不变量。
证明定理 H、I 时，
我们将研究
$HF((L,b);\Lambda_{0,\mathrm{nov}})$
自由部分的秩。

另一方面，
若对两个拉格朗日子流形
分别施加不同的辛微分同胚或哈密顿微分同胚，
则这一对子流形的弗洛尔上同调
在 $\Lambda_{0,\mathrm{nov}}$ 上
不是不变量，
在 $\Lambda_{\mathrm{nov}}$ 上却是不变量。
更准确地说，
$\Lambda_{0,\mathrm{nov}}$ 上
弗洛尔上同调的挠部分并非不变量\liftIndexRangeEnd{I217}。
因此，挠部分可看作描述
对哈密顿同痕的依赖性。
导论中的定理 J
描述了这种现象。

设 $\psi:M\to M$
为哈密顿微分同胚，
$\mu$ 为它到恒等映射的
霍费尔距离\liftIndex{I218}
\textbf{[Hof90, HoZe94]}
（见第 5.3.5 小节）。
假设 $\psi(L^{(0)})$
与 $L^{(1)}$ 横截。
令
\[
b(\mu)=\#\{i\mid\lambda_i\ge\mu\},
\]
其中 $\lambda_i$ 由 (6.1.22) 给出。
则定理 J 陈述不等式
\begin{equation}
\#\bigl(\psi(L^{(0)})\cap L^{(1)}\bigr)
\ge a+2b(\mu).
\tag{6.1.23}\label{eq:6-1-23}
\end{equation}

定理 J 是下面定理 6.1.25 的推论。
后者表明，
弗洛尔上同调的挠指数
相对于霍费尔距离
是利普希茨常数为一的函数
（见第 6.5.3 小节）。
设
$\{\psi_\rho^{(j)}\}_\rho$
（$j=0,1$）为哈密顿同痕，
其中 $\rho\in[0,1]$ 且
$\psi_0^{(j)}$ 为恒等映射。
令 $\psi^{(j)}=\psi_1^{(j)}$，
并记
\begin{equation}
\begin{aligned}
\mu
&=\mathfrak E
\bigl(\{\psi_\rho^{(0)}\}_\rho,
\{\psi_\rho^{(1)}\}_\rho\bigr)\\
&=\bigl\|\{\psi_\rho^{(0)}\}_\rho\bigr\|
+\bigl\|\{\psi_\rho^{(1)}\}_\rho\bigr\|,
\end{aligned}
\tag{6.1.24}\label{eq:6-1-24}
\end{equation}
即 (5.3.40) 中的常数；
其中每个范数表示相应哈密顿同痕的霍费尔长度。
在下一个定理中，
把挠指数 $\lambda_i,\lambda_i'$ 重新编号为
$\lambda_{\downarrow i},
\lambda'_{\downarrow i}$，
使两列分别按非增次序排列。

\begin{theorem}{6.1.25}\label{thm:6-1-25}
设 $\lambda_{\downarrow i}$（$i=1,\ldots,b$）
是
$HF((L^{(1)},\mathbf b_1),
(L^{(0)},\mathbf b_0);
\Lambda_{0,\mathrm{nov}})$
的挠指数，
而 $\lambda'_{\downarrow i}$（$i=1,\ldots,b'$）
是下列弗洛尔上同调的挠指数：
\[
HF\bigl(
(\psi^{(1)}L^{(1)},\psi^{(1)}_*\mathbf b_1),
(\psi^{(0)}L^{(0)},\psi^{(0)}_*\mathbf b_0);
\Lambda_{0,\mathrm{nov}}
\bigr)
\]
假设
$\lambda_{\downarrow i}\ge
\lambda_{\downarrow i+1}$、
$\lambda'_{\downarrow i}\ge
\lambda'_{\downarrow i+1}$，
并令 $\mu$ 如 (6.1.24)。
若 $\lambda_{\downarrow i}>\mu$，
则 $i\le b'$，且
\[
\bigl|\lambda_{\downarrow i}
-\lambda'_{\downarrow i}\bigr|\le\mu.
\]
特别地，
只要 $\lambda_{\downarrow i}>0$，
各 $\lambda_{\downarrow i}$ 连续变化。
\end{theorem}

切卡诺夫在 \textbf{[Chek96]} 中证明：
若 $\psi$ 的霍费尔范数
小于一切边界落在 $L$ 上的
非平凡伪全纯圆盘的面积，
并且 $\psi(L)$ 与 $L$ 横截，
则 $\#(\psi(L)\cap L)$
不小于 $L$ 的贝蒂数。
其结果可由定理 6.5.47 推出；
事实上，我们对该定理的证明
与切卡诺夫的证明相似。
定理 J、6.1.25、6.5.47
可看作利用弗洛尔上同调
对切卡诺夫结果的加强。

我们还将用定理 6.1.25 说明：
对任意相对 spin 对
$(L^{(0)},L^{(1)})$，
即使不一定干净相交，
仍可定义 $\Lambda_{0,\mathrm{nov}}$ 系数的
弗洛尔上同调，
并且它不依赖殆复结构、扰动等
定义时的选择（定义 6.5.39）。

\begin{remark}{6.1.26}\label{rem:6-1-26}
定理 6.1.25 讨论了弗洛尔上同调的
挠指数及其连续性。
挠指数是可以表示为
若干伪全纯圆盘之和的能量的数，
与殆复结构和扰动无关。
它是从集合
\begin{equation}
\bigl\{E(\beta)\mid
\beta\in\pi_2(M,L)
\text{ 可由伪全纯圆盘实现}\bigr\}
\tag{6.1.27}\label{eq:6-1-27}
\end{equation}
中以不依赖各种选择的方式提取出来的。
注意，集合 (6.1.27) 本身
依赖于殆复结构的选择。

利用弗洛尔上同调的乘积结构
$\mathfrak m_k$，
很可能还可以提取更多数值，
使它们不依赖 $J$ 与扰动，
并且相对于霍费尔距离连续。
这些数值可看作
哈密顿微分同胚的谱不变量
在拓扑和拉格朗日情形中的对应物
（见 \textbf{[Oh04]}）。

滤过 $A_\infty$ 代数构造完成，
并证明它在同伦等价意义下
良定义以后，
提取这类信息主要是代数问题，
而非几何或解析问题。
然而，我们还不知道
从滤过 $A_\infty$ 代数或双模
提取数值的最优方法。
因此，这里仅研究挠指数；
它只依赖弗洛尔上同调的加法结构。
\end{remark}

\begin{remark}{6.1.28}\label{rem:6-1-28}
弗洛尔上同调的挠还可能以如下方式
与镜像对称相关。
回忆 $T=q$ 可看作参数
$T\in D^2(\epsilon)$，
这个圆盘参数化镜像中的
极大退化复流形族（见第 1.4 节）。
把 $T$ 求逆，
即用 $\Lambda_{\mathrm{nov}}$
取代 $\Lambda_{0,\mathrm{nov}}$，
对应于用去心圆盘
$D^2(\epsilon)\setminus\{0\}$
取代 $D^2(\epsilon)$；
换言之，这与注 1.23
提到的刚性解析空间有关。
把弗洛尔上同调看作
$\Lambda_{0,\mathrm{nov}}$-模，
而非 $\Lambda_{\mathrm{nov}}$-模，
对应于选定镜像对象
向 $D^2(\epsilon)$ 中零点纤维的延拓。
第 6.5 节的讨论表明，
这种延拓以相当微妙
（而且有趣）的方式
依赖于给定哈密顿同痕类中
代表元的选择。
\end{remark}
